\section{Ordered Set de PBDS}

\whenuse{Conjunto ordenado din\'amico con consultas por orden: obtener el
$k$-\'esimo elemento o contar elementos estrictamente menores que una clave.}

\complexity{insertar, borrar, \texttt{find\_by\_order} y
\texttt{order\_of\_key}: $O(\log n)$; memoria $O(n)$.}

\begin{lstlisting}[style=cpp]
#include <ext/pb_ds/assoc_container.hpp>
#include <ext/pb_ds/tree_policy.hpp>
using namespace __gnu_pbds;

template<class Key>
using OrderedSet = tree<
    Key,
    null_type,
    less<Key>,
    rb_tree_tag,
    tree_order_statistics_node_update
>;

void ordered_set_example() {
    OrderedSet<int> values;

    values.insert(20);
    values.insert(5);
    values.insert(10);

    cout << *values.find_by_order(0) << '\n'; // 5
    cout << *values.find_by_order(1) << '\n'; // 10
    cout << values.order_of_key(15) << '\n';  // 2
    cout << values.order_of_key(10) << '\n';  // 1

    values.erase(10);
}

// Para repetidos, distinguir cada copia con un ID.
using OrderedMultiset = OrderedSet<pair<int,int>>;
\end{lstlisting}

\begin{itemize}
  \item \texttt{find\_by\_order(k)} devuelve un iterador al elemento de
  \'indice $k$, comenzando en cero.
  \item \texttt{order\_of\_key(x)} cuenta los elementos estrictamente menores
  que $x$.
  \item Si $k\ge size()$, \texttt{find\_by\_order(k)} devuelve
  \texttt{end()}; no debe desreferenciarse.
\end{itemize}

\warning{PBDS no pertenece al est\'andar de C++; usar solamente si el juez
compila con GNU G++ y verificar su disponibilidad durante el contest de
pr\'actica.}